\documentclass{article}
\usepackage{amsmath,amssymb}
\usepackage[vocab={planegeometry}]{leanify}
\newtheorem{lemma}{Lemma}
\begin{document}

\section{Notation}

We write $\delta(P,Q)$ for the distance between $P$ and $Q$.
We write $d(P,AB)$ for the distance from $P$ to the line through $A$ and $B$.
We write $H(P,AB)$ for the foot of the perpendicular from $P$ to the line through $A$ and $B$.
We write $O(A,B,C)$ for the circumcentre of the triangle with vertices $A$, $B$ and $C$.
We write $R(A,B,C)$ for the circumradius of the triangle with vertices $A$, $B$ and $C$.
We write $r(A,B,C)$ for the inradius of the triangle with vertices $A$, $B$ and $C$.
We write $E(F,G,a)$ for the ellipse with foci $F$ and $G$ and semi-major axis $a$.
We write $e(F,G,a)$ for the eccentricity of the ellipse with foci $F$ and $G$ and semi-major axis $a$.

\section{Lines and triangles}

\begin{lemma}\label{foot}
Let $A$, $B$ and $P$ be points with $A\ne B$. Then $\delta(P,H(P,AB))=d(P,AB)$.
\end{lemma}

\begin{lemma}\label{circumcentre}
Let $A$, $B$ and $C$ be points. Suppose that $A$ is not collinear with $B$ and $C$. Then $\delta(O(A,B,C),A)=R(A,B,C)$ and $r(A,B,C)\le R(A,B,C)$.
\end{lemma}

\section{Conics}

\begin{lemma}\label{ellipse}
Let $F$ and $G$ be points and let $a$ be a real number with $\delta(F,G)<2a$. Then $e(F,G,a)<1$, and $F$ is not lying on $E(F,G,a)$.
\end{lemma}

\begin{lemma}\label{angle-sum}
Let $A$, $B$ and $C$ be points. Then $|AC|\le|AB|+|BC|$. Suppose that $A$ is not collinear with $B$ and $C$. Then $\angle BAC+\angle ABC+\angle BCA=\pi$.
\end{lemma}

\begin{lemma}\label{non-collinear}
Let $A$, $B$ and $C$ be non-collinear points. Then $A\ne B$ and $\angle BAC<\pi$.
\end{lemma}

\begin{lemma}\label{collinear}
Let $A$, $B$ and $C$ be points that are not collinear. Then $B$ is not collinear with $C$ and $A$. Let $D$, $E$ and $F$ be collinear points. Then $F$ is collinear with $D$ and $E$.
\end{lemma}

\begin{lemma}\label{triangles}
Let $ABC$ be a triangle. Then $\angle A+\angle B+\angle C=\pi$. Let $DEF$ be a triangle with $|DE|=|DF|$. Then $\angle E=\angle F$. Let $PQR$ be a triangle with $\angle R=\pi/2$. Then $|PQ|^2=|PR|^2+|QR|^2$.
\end{lemma}

\end{document}
